\subsection{Логические операции}

Логика: \verb|and| (И), \verb|or| (ИЛИ), \verb|not| (НЕ). Ложными считаются только \verb|false| и \verb|nil| (\hyperlink{term:truthiness}{истинность}): число \verb|0| и пустая строка --- истина. Операторы \verb|and| и \verb|or| вычисляют правую часть только при необходимости («короткое замыкание»).

\begin{codefile}{lua/04_language/lua_logic_ops.lua}
\end{codefile}

\paragraph{Что происходит.}
\begin{itemize}
    \item \verb|gt and 1 or 0| --- идиома «если \verb|gt|, то 1, иначе 0». В Lua нет тернарного оператора, его заменяет эта пара.
    \item \verb|(a >= 0) and (a <= 100)| --- проверка диапазона. Запись \verb|0 <= a <= 100| в Lua не работает.
    \item \verb|name ~= nil and name ~= ""| --- сначала убеждаемся, что значение есть, и только потом сравниваем со строкой. Так защищаются от \verb|nil|.
    \item Строки сравниваются лексикографически: \verb|"a" < "b"| истинно.
\end{itemize}

\subsubsection*{Задачи для самостоятельного решения}
\begin{enumerate}
  \item Для \verb|voltage = 3.4| и \verb|temp = 85| напечатайте «сигнал опасности»: \verb|voltage < 3.5 or temp > 80|.
  \item Напишите предикат «индекс чётный» для \verb|i = 7|: \verb|i % 2 == 0|.
  \item Выберите максимум из \verb|a| и \verb|b| одной строкой через \verb|and/or|.
  \item Инвертируйте булеву переменную \verb|on| и напечатайте её до и после.
  \item Проверьте закон де Моргана: напечатайте \verb|not (p and q)| и \verb|(not p) or (not q)| для всех четырёх сочетаний \verb|p|, \verb|q|.
  \item Напишите условие «\verb|t| не \verb|nil| и у \verb|t| есть поле \verb|value|».
\end{enumerate}
